\section{The $XYZ^2$ code}
\label{sec:code}

\subsection{Matching codes on the honeycomb lattice}

Kitaev's honeycomb model places one qubit on every vertex of a honeycomb lattice and labels every edge $X$, $Y$ or $Z$ according to its direction~\cite{kitaev2006}. An edge $\ell=(j,k)$ with label $\alpha$ carries the link operator
\begin{equation}
  K_\ell=\sigma^{\alpha}_{j}\sigma^{\alpha}_{k},
\end{equation}
and every hexagon $p$ carries the plaquette operator
\begin{equation}
  W_p=\prod_{v\in p}\sigma^{\alpha(v,p)}_{v},
  \label{eq:Wp}
\end{equation}
where $\alpha(v,p)$ is the label of the edge that leaves $p$ at the vertex $v$. Every plaquette commutes with every link operator, while link operators on edges that share a vertex anticommute. A matching code takes as stabilizer generators all plaquettes together with the link operators of a matching, that is, of a set of edges with no common vertex~\cite{wootton}. These codes have the anyon content of the surface code, and their syndromes can be decoded by minimum-weight matching~\cite{wootton,wootton2022}. Subsystem codes on the same lattice measure only two-body link operators~\cite{suchara2011,wootton2022}.

\subsection{Definition of the code}

The $XYZ^2$ code is the matching code in which the matching consists of the vertical edges, which carry the label $X$~\cite{Sri2022}. We call each vertical edge a block and label the blocks of a distance-$d$ patch by $(i,j)$ with $0\le i,j\le d-1$. Block $(i,j)$ has an upper vertex $u_{ij}$ and a lower vertex $l_{ij}$, so the patch has $n=2d^2$ qubits. \autoref{fig:code}(a) shows the patch for $d=3$. The stabilizer group is generated by the following operators:
\begin{enumerate}
  \item the $d^2$ links $L_{ij}=X_{u_{ij}}X_{l_{ij}}$,
  \item the $(d-1)^2$ hexagons
  \begin{equation}
    H_{ij}=X_{l_{i,j}}\,Y_{u_{i,j+1}}\,Z_{l_{i,j+1}}\,X_{u_{i+1,j+1}}\,Y_{l_{i+1,j}}\,Z_{u_{i+1,j}}
    \label{eq:hexagon}
  \end{equation}
  for $0\le i,j\le d-2$, which follow \autoref{eq:Wp},
  \item $2(d-1)$ boundary checks of weight three.
\end{enumerate}
A boundary check is the hexagon of a cell $(i,j)$ just outside the patch, restricted to its three qubits inside the patch. On the upper-right and lower-left edges of the patch we use the cells with $i+j$ even, and on the upper-left and lower-right edges the cells with $i+j$ odd [\autoref{fig:code}(a)]. These $2d^2-1$ generators are independent, so the code encodes one logical qubit. Its distance is $d$~\cite{Sri2022}.

\begin{figure*}[t]
  \centering
  \includegraphics[width=0.86\textwidth]{fig_code.pdf}
  \caption{The distance-3 $XYZ^2$ code. (a)~The 18 qubits sit on the vertices of a honeycomb patch. Thick vertical edges are the nine $XX$ links. Block $(0,0)$ is the top link, with upper vertex $u_{00}$ and lower vertex $l_{00}$, and the indices $i$ and $j$ grow toward the lower left and the lower right. The four shaded hexagons and the four dashed triangles are the weight-six and weight-three plaquette checks. The Pauli labels on hexagon $H_{10}$ show the pattern $XYZXYZ$ of \autoref{eq:hexagon}. Coral cells have $i+j$ even (class B) and sunflower cells have $i+j$ odd (class A). (b)~Each link is a $[[2,1,1]]$ code that holds one effective qubit. On the $3\times3$ array of effective qubits every plaquette acts as $\eff{Z}$ on two diagonally opposite blocks and as $\eff{Y}$ on the other two, which defines the YZZY surface code. (c)~The block-local Clifford map of \autoref{sec:frame} turns the YZZY code into a CSS surface code, with $\eff{Z}$-type class-B plaquettes and $\eff{X}$-type class-A plaquettes. The logical operator $\Xb$ (violet band) becomes a $\eff{Z}$-type string, and $\plusL$ becomes the logical zero state, which is prepared from effective zero states.}
  \label{fig:code}
\end{figure*}

The logical operator
\begin{equation}
  \Xb=\prod_{i+j=d-1}X_{l_{ij}}
  \label{eq:xbar}
\end{equation}
acts on the lower vertices of the $d$ blocks on the anti-diagonal $i+j=d-1$, which form the middle row of the patch. It has weight $d$ and contains only $X$ operators. The operators $\Yb$ and $\Zb$ mix Pauli types~\cite{Sri2022}. \autoref{fig:logicals} shows the three logical operators for $d=3$. In this paper we prepare $\plusL$, the $+1$ eigenstate of $\Xb$. Only errors that anticommute with $\Xb$ can corrupt this state. For the distances studied here, the lightest logical operator that anticommutes with $\Xb$ has weight $d+1$, as $\Yb$ does in \autoref{fig:logicals}(b). The preparation of $\plusL$ is therefore protected to distance $d+1$, one more than the code distance.

\begin{figure}[t]
  \centering
  \includegraphics[width=\columnwidth]{fig_logicals.pdf}
  \caption{Logical operators of the $d=3$ code. (a)~$\Xb=X_9X_{10}X_{11}$ has weight 3. (b)~$\Yb=X_4Z_7Y_{10}X_{13}$ has weight 4. (c)~$\Zb=X_4Z_7X_9Z_{10}X_{11}X_{13}$ is the product of $\Xb$ and $\Yb$ up to a phase.}
  \label{fig:logicals}
\end{figure}

\subsection{Block structure}
\label{sec:blocks}

Each block carries a $[[2,1,1]]$ phase-flip code with stabilizer $L_{ij}$. We take
\begin{equation}
  \eff{Z}_{ij}=X_{u_{ij}},\qquad \eff{X}_{ij}=Z_{u_{ij}}Z_{l_{ij}}
\end{equation}
as its logical operators, so that $\eff{Y}_{ij}\propto Y_{u_{ij}}Z_{l_{ij}}$, which equals $Z_{u_{ij}}Y_{l_{ij}}$ up to a sign and a factor of $L_{ij}$. Multiplying a hexagon by links does not change the code. Up to such factors, every hexagon acts on its four blocks as
\begin{equation}
  H_{ij}\simeq \eff{Z}_{i,j}\,\eff{Y}_{i,j+1}\,\eff{Z}_{i+1,j+1}\,\eff{Y}_{i+1,j}.
  \label{eq:yzzy}
\end{equation}
The blocks therefore carry a $d\times d$ surface code with YZZY plaquettes, and the boundary checks become its weight-two boundary stabilizers [\autoref{fig:code}(b)]. The $XYZ^2$ code is this YZZY surface code concatenated with the $d^2$ link codes~\cite{Sri2022,Sri2025}.

We call block $(i,j)$ even if $i+j$ is even and odd otherwise. The same parity splits the cells into two classes. A class-B cell ($i+j$ even) acts as $\eff{Z}$ on two even blocks and as $\eff{Y}$ on two odd blocks. A class-A cell ($i+j$ odd) acts as $\eff{Z}$ on two odd blocks and as $\eff{Y}$ on two even blocks. The boundary checks on the upper-right and lower-left edges belong to class B and those on the upper-left and lower-right edges to class A. This bipartition plays the role of the $X$ and $Z$ plaquettes of the CSS surface code in \autoref{sec:protocol}.

The concatenated structure also explains how errors appear in the syndrome. A $Z$ error on either vertex of block $(i,j)$ flips $L_{ij}$. The two errors differ by $\eff{X}_{ij}$, so the link syndrome reveals that the block was hit but not which vertex. The sequential decoder of Ref.~\cite{Sri2025} removes every link syndrome with a fixed choice of vertex and passes the resulting conditional error probabilities of the effective qubits to a decoder of the YZZY code. We return to this two-level structure in \autoref{sec:decoding}.
